\documentclass[UTF8,11pt]{ctexart}
\usepackage[a4paper,margin=22mm]{geometry}
\usepackage{graphicx,enumitem,amssymb,tabularx}
\usepackage[normalem]{ulem}
\usepackage{xeCJKfntef}
\setlength{\parindent}{0pt}
\setlength{\parskip}{.6em}
\sloppy
\begin{document}
\section*{2016年6月英语四级真题(第1套)}
\textbf{2016} 年\textbf{6} 月大学英语四级真题（第\textbf{1} 套）


\subsection*{\textbf{Part I} \textbf{Writing} \textbf{(30 minutes)}}


Directions: For this part, you are allowed 30 minutes to write a letter to\textbf{ express your thanks to your parents or any} \textbf{family members upon making memorable achievement.} You should write at least 120 words but no more than 180 words.


\subsection*{\textbf{Part} II \textbf{Listening Comprehension} \textbf{(25 minutes)}}


\subsection*{\textbf{Section A}}


\textbf{Directions:} In this section, you will hear three news reports. At the end of each news report, you will hear two or three questions. Both the news report and the questions will be spoken only once. After you hear a question, you must choose the best answer from the four choices marked A), B), C) and D). Then mark the corresponding letter on Answer Sheet 1 with a single line through the centre.


\textbf{Questions 1 and 2 are based on the news report you have just heard.}


1.


\begin{itemize}[leftmargin=2em,itemsep=.2em,topsep=.4em]\item[A.] The International Labor Organization's key objective.
\item[B.] The basic social protection for the most vulnerable.
\item[C.] Rising unemployment worldwide.
\item[D.] Global economic recovery.
\end{itemize}

2.


\begin{itemize}[leftmargin=2em,itemsep=.2em,topsep=.4em]\item[A.] Many countries have not taken measures to create enough jobs.
\item[B.] Few countries know how to address the current economic crisis.
\item[C.] Few countries have realized the seriousness of the current crisis.
\item[D.] Many countries need support to improve their people's livelihood.
\end{itemize}

Questions 3 and 4 are based on the news report you have just heard.


\textbf{Questions 3and 4 are based on the news report you have just heard.}


3.


\begin{itemize}[leftmargin=2em,itemsep=.2em,topsep=.4em]\item[A.] Serve standardized food nationwide.
\item[B.] Put calorie information on the menu.
\item[C.] Increase protein content in the food.
\item[D.] Offer convenient food to customers.
\end{itemize}

\noindent\begin{tabularx}{\linewidth}{@{}lXXXX@{}}4. & \textbf{A.} They will be fined. & \textbf{B.} They will be closed. & \textbf{C.} They will get a warning. & \textbf{D.} They will lose customers.\\\end{tabularx}\par

\textbf{Questions 5 to 7 are based on the news report you have just heard.}


5.


\begin{itemize}[leftmargin=2em,itemsep=.2em,topsep=.4em]\item[A.] Inability to implement their business plans.
\item[B.] Inability to keep turning out novel products.
\item[C.] Lack of a successful business model of their own.
\item[D.] Failure to integrate innovation into their business.
\end{itemize}

6.


\begin{itemize}[leftmargin=2em,itemsep=.2em,topsep=.4em]\item[A.] It is the secret to business success.
\item[B.] It is the creation of something new.
\item[C.] It is a magic tool to bring big rewards.
\item[D.] It is an essential part of business culture.
\end{itemize}

7.


\begin{itemize}[leftmargin=2em,itemsep=.2em,topsep=.4em]\item[A.] Its hardworking employees.
\item[B.] Its flexible promotion strategy.
\item[C.] Its innovation culture.
\item[D.] Its willingness to make investments.
\end{itemize}

\subsection*{\textbf{Section B}}


\textbf{Directions:} In this section, you will hear two long conversations. At the end of each conversations you will hear four questions. Both the conversations and the question-s will be spoken only once. After you hear a question. You must choose the best answer from the four choices marked A), B), C)and D). Then mark the corresponding letter on Answer Sheet 1 with a single line through the centre.


\textbf{Questions 8 to 11 are based on the conversation you have just heard.}


8.


\begin{itemize}[leftmargin=2em,itemsep=.2em,topsep=.4em]\item[A.] He's got addicted to technology.
\item[B.] He is not very good at socializing.
\item[C.] He is crazy about text-messaging.
\item[D.] He does not talk long on the phone.
\end{itemize}

\noindent\begin{tabularx}{\linewidth}{@{}lXXXX@{}}9. & \textbf{A.} Talk big. & \textbf{B.} Talk at length. & \textbf{C.} Gossip a lot. & \textbf{D.} Forget herself.\\\end{tabularx}\par

10.


\begin{itemize}[leftmargin=2em,itemsep=.2em,topsep=.4em]\item[A.] He thought it was cool.
\item[B.] He needed the practice.
\item[C.] He wanted to stay connected with them.
\item[D.] He had an urgent message to send.
\end{itemize}

11.


\begin{itemize}[leftmargin=2em,itemsep=.2em,topsep=.4em]\item[A.] It poses a challenge to seniors.
\item[B.] It saves both time and money.
\item[C.] It is childish and unprofessional.
\item[D.] It is cool and convenient.
\end{itemize}

\textbf{Questions 12 to 15 are based on the conversation you have just heard.}


12.


\begin{itemize}[leftmargin=2em,itemsep=.2em,topsep=.4em]\item[A.] He wants to change his job assignment.
\item[B.] He is unhappy with his department manager.
\item[C.] He thinks he deserves extra pay for overtime.
\item[D.] He is often singled out for criticism by his boss.
\end{itemize}

13.


\begin{itemize}[leftmargin=2em,itemsep=.2em,topsep=.4em]\item[A.] His workload was much too heavy.
\item[B.] His immediate boss did not trust him.
\item[C.] His colleagues often refused to cooperate.
\item[D.] His salary was too low for his responsibility.
\end{itemize}

14.


\begin{itemize}[leftmargin=2em,itemsep=.2em,topsep=.4em]\item[A.] He never knows how to refuse.
\item[B.] He is always ready to help others.
\item[C.] His boss has a lot of trust in him.
\item[D.] His boss has no sense of fairness.
\end{itemize}

15.


\begin{itemize}[leftmargin=2em,itemsep=.2em,topsep=.4em]\item[A.] Put all his complaints in writing.
\item[B.] Wait and see what happens next.
\item[C.] Learn to say no when necessary.
\item[D.] Talk to his boss in person first.
\end{itemize}

\subsection*{\textbf{Section C}}


\textbf{Directions:} In this section, you will hear three passages. At the end of each passage, you will hear three or four questions. Both the passage and the questions will be spoken only once. After you hear a question, you must choose the best answer from the four choices marked A), B), C)and D).Then mark the corresponding letter on Answer Sheet 1 with a single line through the centre.


\textbf{Questions 16 to 18 are based on the passage you have just heard.}


16.


\begin{itemize}[leftmargin=2em,itemsep=.2em,topsep=.4em]\item[A.] The importance of sleep to a healthy life.
\item[B.] Reasons for Americans' decline in sleep.
\item[C.] Some tips to improve the quality of sleep.
\item[D.] Diseases associated with lack of sleep.
\end{itemize}

17.


\begin{itemize}[leftmargin=2em,itemsep=.2em,topsep=.4em]\item[A.] They are more health-conscious.
\item[B.] They are changing their living habits.
\item[C.] They get less and less sleep.
\item[D.] They know the dangers of lack of sleep.
\end{itemize}

18.


\begin{itemize}[leftmargin=2em,itemsep=.2em,topsep=.4em]\item[A.] Their weight will go down.
\item[B.] Their mind function will deteriorate.
\item[C.] Their work efficiency will decrease.
\item[D.] Their blood pressure will rise.
\end{itemize}

\textbf{Questions 19 to 21 are based on the passage you have just heard.}


19.


\begin{itemize}[leftmargin=2em,itemsep=.2em,topsep=.4em]\item[A.] How much you can afford to pay.
\item[B.] What course you are going to choose.
\item[C.] Which university you are going to apply to.
\item[D.] When you are going to submit your application.
\end{itemize}

20.


\begin{itemize}[leftmargin=2em,itemsep=.2em,topsep=.4em]\item[A.] The list of courses studied.
\item[B.] The full record of scores.
\item[C.] The references from teachers.
\item[D.] The personal statement.
\end{itemize}

21.


\begin{itemize}[leftmargin=2em,itemsep=.2em,topsep=.4em]\item[A.] Specify what they would like to do after graduation.
\item[B.] Describe in detail how much they would enjoy studying.
\item[C.] Indicate they have reflected and thought about the subject.
\item[D.] Emphasize that they admire the professors in the university.
\end{itemize}

\textbf{Questions 22 to 25 are based on the passage you have just heard.}


22.


\begin{itemize}[leftmargin=2em,itemsep=.2em,topsep=.4em]\item[A.] It was equipped with rubber tyres.
\item[B.] It was built in the late 19th century.
\item[C.] It was purchased by the Royal family.
\item[D.] It was designed by an English engineer.
\end{itemize}

23.


\begin{itemize}[leftmargin=2em,itemsep=.2em,topsep=.4em]\item[A.] They consumed lots of petrol.
\item[B.] They took two passengers only.
\item[C.] They were difficult to drive.
\item[D.] They often broke down.
\end{itemize}

24.


\begin{itemize}[leftmargin=2em,itemsep=.2em,topsep=.4em]\item[A.] They were produced on the assembly line.
\item[B.] They were built with less costly materials.
\item[C.] They were modeled after British cars.
\item[D.] They were made for ordinary use.
\end{itemize}

25.


\begin{itemize}[leftmargin=2em,itemsep=.2em,topsep=.4em]\item[A.] It made news all over the world.
\item[B.] It was built for the Royal family.
\item[C.] It marked a new era in motor travel.
\item[D.] It attracted large numbers of motorists.
\end{itemize}

\subsection*{\textbf{Part III} \textbf{Reading Comprehension}}


\subsection*{\textbf{Section A}}


\textbf{Directions:}\textit{ In this section, there is a passage with ten blanks. You are required to select one word for each blank} \textit{from a list of choices given in a word bank following the passage. Read the passage through carefully before making} \textit{your choices. Each choice in the bank is identified by a letter. Please mark the corresponding letter for each item on} \textbf{Answer Sheet 2}\textit{ with a single line through the centre. You may not use any of the words in the bank more than once.}


Physical activity does the body good, and there's growing evidence that it helps the brain too. Researchers in the Netherlands report that children who get more exercise, whether at school or on their own, \uline{26} to have higher GPAs and better scores on standardized tests. In a \uline{27} of 14 studies that looked at physical activity and academic \uline{28} , investigators found that the more children moved, the better their grades were in school, \uline{29} in the basic subjects of math, English and reading.


The data will certainly fuel the ongoing debate over whether physical education classes should be cut as schools struggle to \uline{30} on smaller budgets. The arguments against physical education have included concerns that gym time may be taking away from study time. With standardized test scores in the U.S. \uline{31} in recent years, some administrators believe students need to spend more time in the classroom instead of on the playground. But as these findings show, exercise and academics may not be \uline{32} exclusive. Physical activity can improve blood \uline{33} to the brain, fueling memory, attention and creativity, which are \uline{34} to learning. And exercise releases hormones that can improve \uline{35} and relieve stress, which can also help learning. So while it may seem as if kids are just exercising their bodies when they're running around, they may actually be exercising their brains as well.


\begin{itemize}[leftmargin=2em,itemsep=.2em,topsep=.4em]\item[A.] attendance
\item[B.] consequently
\item[C.] current
\item[D.] depressing
\item[E.] dropping
\item[F.] essential
\item[G.] feasible
\item[H.] flow
\item[I.] mood
\item[J.] mutually
\item[K.] particularly
\item[L.] performance
\item[M.] review
\item[N.] survive
\item[O.] tend
\end{itemize}

\subsection*{\textbf{Section B}}


\textbf{Directions:}\textit{ In this section, you are going to read a passage with ten statements attached to it. Each statement} \textit{contains information given in one of the paragraphs. Identify the paragraph from which the information is derived.} \textit{You may choose a paragraph more than once. Each paragraph is marked with a letter. Answer the questions by} \textit{marking the corresponding letter on}\textbf{ Answer Sheet 2}.


\subsection*{\textbf{Finding the Right Home—and Contentment, Too}}


\begin{itemize}[leftmargin=2em,itemsep=.2em,topsep=.4em]\item[A.] When your elderly relative needs to enter some sort of long-term care facility—a moment few parents or children
\end{itemize}

approach without fear—what you would like is to have everything made clear.


\begin{itemize}[leftmargin=2em,itemsep=.2em,topsep=.4em]\item[B.] Does assisted living really mark a great improvement over a nursing home, or has the industry simply hired better
\end{itemize}

interior designers? Are nursing homes as bad as people fear, or is that an out-moded\textit{ stereotype}（固定看法)? Can doing one's homework really steer families to the best places? It is genuinely hard to know.


\begin{itemize}[leftmargin=2em,itemsep=.2em,topsep=.4em]\item[C.] I am about to make things more complicated by suggesting that what kind of facility an older person lives in
\end{itemize}

may matter less than we have assumed. And that the characteristics adult children look for when they begin the search are not necessarily the things that make a difference to the people who are going to move in. I am not talking about the quality of care, let me hastily add. Nobody flourishes in a gloomy environment with irresponsible staff and a poor safety record. But an accumulating body of research indicates that some distinctions between one type of elder care and another have little real bearing on how well residents do.


\begin{itemize}[leftmargin=2em,itemsep=.2em,topsep=.4em]\item[D.] The most recent of these studies, published in\textit{ The journal of Applied Gerontology}, surveyed 150 Connecticut
\end{itemize}

residents of assisted living, nursing homes and smaller residential care homes （known in some states as board and care homes or adult care homes). Researchers from the University of Connecticut Health Center asked the residents a large number of questions about their quality of life, emotional well-being and social interaction, as well as about the quality of the facilities.


\begin{itemize}[leftmargin=2em,itemsep=.2em,topsep=.4em]\item[E.] “We thought we would see differences based on the housing types,” said the lead author of the study,Julie
\end{itemize}

Robison, an associate professor of medicine at the university. A reasonable assumption—don't families struggle to avoid nursing homes and suffer real guilt if they can't?


\begin{itemize}[leftmargin=2em,itemsep=.2em,topsep=.4em]\item[F.] In the initial results, assisted living residents did paint the most positive picture. They were less likely to report
\end{itemize}

symptoms of depression than those in the other facilities, for instance, and less likely to be bored or lonely. They scored higher on social interaction.


\begin{itemize}[leftmargin=2em,itemsep=.2em,topsep=.4em]\item[G.] But when the researchers plugged in a number of other variables, such differences disappeared. It is not the
\end{itemize}

housing type, they found, that creates differences in residents' responses. “It is the characteristics of the specific environment they are in, combined with their own personal characteristics—how healthy they feel they are, their age and marital status,” Dr. Robison explained. Whether residents felt involved in the decision to move and how long they had lived there also proved significant.


\begin{itemize}[leftmargin=2em,itemsep=.2em,topsep=.4em]\item[H.] An elderly person who describes herself as in poor health, therefore, might be no less depressed in assisted living
\end{itemize}

（even if her children preferred it) than in a nursing home. A person who had input into where he would move and has had time to adapt to it might do as well in a nursing home as in a small residential care home, other factors being equal. It is an interaction between the person and the place, not the sort of place in itself, that leads to better or worse experiences.“You can't just say,'Let's put this person in a residential care home instead of a nursing home—she will be much better off,'” Dr. Robison said. What matters, she added, “is a combination of what people bring in with them, and what they find there.”


\begin{itemize}[leftmargin=2em,itemsep=.2em,topsep=.4em]\item[I.] Such findings, which run counter to common sense, have surfaced before. In a multi-state study of assisted living,
\end{itemize}

for instance, University of North Carolina researchers found that a host of variables—the facility's type, size or age； whether a chain owned it；how attractive the neighborhood was—had no significant relationship to how the residents fared in terms of illness, mental decline, hospitalizations or mortality. What mattered most was the residents' physical health and mental status. What people were like when they came in had greater consequence than what happened once they were there.


\begin{itemize}[leftmargin=2em,itemsep=.2em,topsep=.4em]\item[J.] As I was considering all this, a press release from a respected research firm crossed my desk, announcing that the
\end{itemize}

five-star rating system that Medicare developed in 2008 to help families compare nursing home quality also has little relationship to how satisfied its residents or their family members are. As a matter of fact, consumers expressed higher satisfaction with the one-star facilities, the lowest rated, than with the five-star ones. （More on this study and the star ratings will appear in a subsequent post.)


\begin{itemize}[leftmargin=2em,itemsep=.2em,topsep=.4em]\item[K.] Before we collectively tear our hair out—how are we supposed to find our way in a landscape this
\end{itemize}

confusing?—here is a thought from Dr. Philip Sloane, a\textit{ geriatrician}（老年病学专家)at the University of North Carolina :“In a way, that could be liberating for families.”


\begin{itemize}[leftmargin=2em,itemsep=.2em,topsep=.4em]\item[L.] Of course, sons and daughters want to visit the facilities, talk to the administrators and residents and other
\end{itemize}

families, and do everything possible to fulfill their duties. But perhaps they don't have to turn themselves into private investigators or Congressional subcommittees. “Families can look a bit more for where the residents are going to be happy,” Dr. Sloane said. And involving the future resident in the process can be very important.


\begin{itemize}[leftmargin=2em,itemsep=.2em,topsep=.4em]\item[M.] We all have our own ideas about what would bring our parents happiness. They have their ideas, too. A friend
\end{itemize}

recently took her mother to visit an expensive assisted living/nursing home near my town. I have seen this place—it is elegant, inside and out. But nobody greeted the daughter and mother when they arrived, though the visit had been planned；nobody introduced them to the other residents. When they had lunch in the dining room, they sat alone at a table.


\begin{itemize}[leftmargin=2em,itemsep=.2em,topsep=.4em]\item[N.] The daughter feared her mother would be ignored there, and so she decided to move her into a more welcoming
\end{itemize}

facility. Based on what is emerging from some of this research, that might have been as rational a way as any to reach a decision.


36. Many people feel guilty when they cannot find a place other than a nursing home for their parents.


37.Though it helps for children to investigate care facilities, involving their parents in the decision-making process may prove very important.


38.It is really difficult to tell if assisted living is better than a nursing home.


39.How a resident feels depends on an interaction between themselves and the care facility they live in.


40.The author thinks her friend made a rational decision in choosing a more hospitable place over an apparently elegant assisted living home.


41.The system Medicare developed to rate nursing home quality is of little help to finding a satisfactory place.


42.At first the researchers of the most recent study found residents in assisted living facilities gave higher scores on social interaction.


43.What kind of care facility old people live in may be less important than we think.


44.The findings of the latest research were similar to an earlier multi-state study of assisted living.


45.A resident's satisfaction with a care facility has much to do with whether they had participated in the decision to move in and how long they had stayed there.


\subsection*{\textbf{Section C}}


\textbf{Directions:}\textit{ There are 2 passages in this section. Each passage is followed by some questions or unfinished} \textit{statements. For each of them there are four choices marked A), B), C) and D). You should decide on the best choice} \textit{and mark the corresponding letter on}\textbf{ Answer Sheet 2}\textit{ with a single line through the centre.}


\subsection*{\textbf{Passage One}}


\textbf{Questions 46 to 50 are based on the following passage.}


As Artificial Intelligence (AI) becomes increasingly sophisticated, there are growing concerns that robots could become a threat. This danger can be avoided, according to computer science professor Stuart Russell, if we figure out how to turn human values into a programmable code.


Russell argues that as robots take on more complicated tasks, it's necessary to translate our morals into AI language.


For example, if a robot does chores around the house, you wouldn't want it to put the pet cat in the oven to make dinner for the hungry children. “You would want that robot preloaded with a good set of values,” said Russell.


Some robots are already programmed with basic human values. For example, mobile robots have been programmed to keep a comfortable distance from humans. Obviously there are cultural differences, but if you were talking to another person and they came up close in your personal space, you wouldn't think that's the kind of thing a properly brought-up person would do.


It will be possible to create more sophisticated moral machines, if only we can find a way to set out human values as clear rules.


Robots could also learn values from drawing patterns from large sets of data on human behavior. They are dangerous only if programmers are careless.


The biggest concern with robots going against human values is that human beings fail to do sufficient testing and they've produced a system that will break some kind of\textit{ taboo}（禁忌).


One simple check would be to program a robot to check the correct course of action with a human when presented with an unusual situation.


If the robot is unsure whether an animal is suitable for the microwave, it has the opportunity to stop, send out \textit{beeps}（嘟嘟声), and ask for directions from a human. If we humans aren't quite sure about a decision, we go and ask somebody else.


The most difficult step in programming values will be deciding exactly what we believe in moral, and how to create a set of ethical rules. But if we come up with an answer, robots could be good for humanity.


46.What does the author say about the threat of robots?


\begin{itemize}[leftmargin=2em,itemsep=.2em,topsep=.4em]\item[A.] It may constitute a challenge to computer programmers.
\item[B.] It accompanies all machinery involving high technology.
\end{itemize}

\begin{itemize}[leftmargin=2em,itemsep=.2em,topsep=.4em]\item[C.] It can be avoided if human values are translated into their language.
\item[D.] It has become an inevitable peril as technology gets more sophisticated.
\end{itemize}

47.What would we think of a person who invades our personal space according to the author?


\begin{itemize}[leftmargin=2em,itemsep=.2em,topsep=.4em]\item[A.] They are aggressive.
\item[B.] They are outgoing.
\item[C.] They are ignorant.
\item[D.] They are ill-bred.
\end{itemize}

48.How do robots learn human values?


\begin{itemize}[leftmargin=2em,itemsep=.2em,topsep=.4em]\item[A.] By interacting with humans in everyday life situations.
\item[B.] By following the daily routines of civilized human beings.
\item[C.] By picking up patterns from massive data on human behavior.
\item[D.] By imitating the behavior of property brought-up human beings.
\end{itemize}

49.What will a well-programmed robot do when facing an unusual situation?


\begin{itemize}[leftmargin=2em,itemsep=.2em,topsep=.4em]\item[A.] Keep a distance from possible dangers.
\item[B.] Stop to seek advice from a human being.
\item[C.] Trigger its built-in alarm system at once.
\item[D.] Do sufficient testing before taking action.
\end{itemize}

50.What is most difficult to do when we turn human values into a programmable code?


\begin{itemize}[leftmargin=2em,itemsep=.2em,topsep=.4em]\item[A.] Determine what is moral and ethical.
\item[B.] Design some large-scale experiments.
\item[C.] Set rules for man-machine interaction.
\item[D.] Develop a more sophisticated program.
\end{itemize}

\subsection*{\textbf{Passage Two}}


\textbf{Questions 51 to 55 are based on the following passage.}


Why do some people live to be older than others? You know the standard explanations: keeping a moderate diet, engaging in regular exercise, etc. But what effect does your personality have on your\textit{ longevity}（长寿）?Do some kinds of personalities lead to longer lives? A new study in the\textit{ Journal of the American Geriatrics Society} looked at this question by examining the personality characteristics of 246 children of people who had lived to be at least 100.


The study shows that those living the longest are more outgoing, more active and less\textit{ neurotic}（神经质的) than other people. Long-living women are also more likely to be sympathetic and cooperative than women with a normal life span. These findings are in agreement with what you would expect from the evolutionary theory: those who like to make friends and help others can gather enough resources to make it through tough times.


Interestingly, however, other characteristics that you might consider advantageous had no impact on whether study participants were likely to live longer. Those who were more self-disciplined, for instance, were no more likely to live to be very old. Also, being open to new ideas had no relationship to long life, which might explain all those bad-tempered old people who are fixed in their ways.


Whether you can successfully change your personality as an adult is the subject of a longstanding psychological debate. But the new paper suggests that if you want long life, you should strive to be as outgoing as possible.


Unfortunately, another recent study shows that your mother's personality may also help determine your longevity. That study looked at nearly 28,000 Norwegian mothers and found that those moms who were more anxious, depressed and angry were more likely to feed their kids unhealthy diets. Patterns of childhood eating can be hard to break when we're adults, which may mean that kids of depressed moms end up dying younger.


Personality isn't\textit{ destiny}（命运), and everyone knows that individuals can learn to change. But both studies show that long life isn't just a matter of your physical health but of your mental health.


51. The aim of the study in the\textit{ Journal of the American Geriatrics Society} is\_\_\_\_.


\begin{itemize}[leftmargin=2em,itemsep=.2em,topsep=.4em]\item[A.] to see whether people's personality affects their life span
\item[B.] to find out if one's lifestyle has any effect on their health
\item[C.] to investigate the role of exercise in living a long life
\item[D.] to examine all the factors contributing to longevity
\end{itemize}

52. What does the author imply about outgoing and sympathetic people?


\begin{itemize}[leftmargin=2em,itemsep=.2em,topsep=.4em]\item[A.] They have a good understanding of evolution.
\item[B.] They are better at negotiating an agreement.
\end{itemize}

\begin{itemize}[leftmargin=2em,itemsep=.2em,topsep=.4em]\item[C.] They generally appear more resourceful.
\item[D.] They are more likely to get over hardship.
\end{itemize}

53. What finding of the study might prove somewhat out of our expectation?


\begin{itemize}[leftmargin=2em,itemsep=.2em,topsep=.4em]\item[A.] Easy-going people can also live a relatively long life.
\item[B.] Personality characteristics that prove advantageous actually vary with times.
\item[C.] Such personality characteristics as self-discipline have no effect on longevity.
\item[D.] Readiness to accept new ideas helps one enjoy longevity.
\end{itemize}

54. What does the recent study of Norwegian mothers show?


\begin{itemize}[leftmargin=2em,itemsep=.2em,topsep=.4em]\item[A.] Children's personality characteristics are invariably determined by their mothers.
\item[B.] People with unhealthy eating habits are likely to die sooner.
\item[C.] Mothers' influence on children may last longer than fathers'.
\item[D.] Mothers' negative personality characteristics may affect their children's life spans.
\end{itemize}

55.What can we learn from the findings of the two new studies?


\begin{itemize}[leftmargin=2em,itemsep=.2em,topsep=.4em]\item[A.] Anxiety and depression more often than not cut short one's life span.
\item[B.] Longevity results from a combination of mental and physical health.
\item[C.] Personality plays a decisive role in how healthy one is.
\item[D.] Health is in large part related to one's lifestyle.
\end{itemize}

\subsection*{\textbf{Part IV} \textbf{Translation}}


\textbf{Directions:}\textit{ For this part, you are allowed 30 minutes to translate a passage from Chinese into English. You should} \textit{write your answer on}\textbf{ Answer Sheet 2}\textit{.}


功夫（Kong Fu) 是中国武术（martial arts）的俗称。中国武术的起源可以追溯到自卫的需要，狩猎活动以及古代中国的军士训练。它是中国传统体育运动的一种，年轻人和老年人都练。它已逐渐演变成了中国文化的独特元素。作为中国的国宝，武术有上百种不同的风格，是世界上练得最多的武术形式。有些风格模仿了动物的动作，还有一些则受到了中国哲学思想，神话和传说的启发。

\end{document}